\subsubsection*{ Fully Autonomous Focused Exploration for Robotics Environmental Monitoring \cite{Hitz2001} } 

Environmental phenomena can sometimes require large datasets to be collected in order to more completely encompass the temporal and spatial dynamics. 

There is a niche for inexpensive robotic platforms that can use sensors to measure along a vertical water column. Toxic floral species, such as Cyanobacteria are often the target for environmental monitoring efforts due to their pervasive and ubiquitous nature. 


\subsubsection*{ Autonomous Inland Water Monitoring, \cite{Hitz2012} } 

\subsubsection*{ Robotics for Environmental Monitoring \cite{Dunbabin2012} } 

\subsubsection*{ An Autonomous Surface Vehicle for Water Quality Monitoring \cite{Dunbabin2009} } 

\subsubsection*{ A Robotic System to Locate Hazardous Chemical Leaks \cite{Russell2014}} 

\subsubsection*{ Bacterium-inspired Robots for Environmental Monitoring \cite{Dhariwal2004} }


\cite{Seto2016} due to volumetric and something constraints placed upon small AUVs. Unique solutions are required in order to effectively take measurements from a given location. Some strategies such as show by \cite{Seto2016} employ a novel methodology that attempts to reduce the number of times that subsurface craft have to surface for GPS locations, greatly speeding up the measuremment process. In this method, limited acoustic feedback is broadcast between the individual robots and then used to more accurately precise their location in space. 